\section{Researcher Control and Experimentability}
\label{sec:eval-control}

The fourth research question concerns the control and the auditable record the researcher needs to operate experiments.
The recorded sessions were operated entirely through the researcher console of \cref{fig:impl-screens-console}: the session could not start until the setup prerequisites and calibration had passed, gated through the stepper of \cref{fig:impl-setup}, the interventions in \cref{tab:eval-sessions} were applied from the console during the run, and the participant's view was mirrored to the researcher throughout.
We verified these flows by walkthrough, as recorded in \cref{tab:req-coverage}.

The auditable record is demonstrated by this chapter itself.
Every figure and table here was produced from the exported session records alone, without reference to any live state, which shows that a session can be reconstructed and analysed from its record (NFR5).
The export carries the session parameters, the calibration metrics, the presentation configuration, the gaze stream, the derived events, and the intervention and context-preservation outcomes, which is what allowed the sensing, latency, and context-preservation results above to be computed after the fact.
Session replay reconstructs a recorded session in the reading interface from its exported record alone, shown in \cref{fig:impl-replay}, which is the same reproducibility property that this chapter exercises throughout (FR21).
